%% kapa: exam=2 round=30 year=2019 period=3 subject=compensation_law subject_ko=감정평가 및 보상법규 seq=3 number=3 points=20
%% source: Q-Net 감정평가사 기출문제 (cst00309) artlSeq=5205595 file=raw/1008/5205595/감정평가 및 보상법규.pdf
%% conversion: kordoc 4.24.51 -> tools/gen2.py (auto); review=proofread-gpt-terra; flags=none
\begin{longquestion}{【문제 3】}
X군에 거주하는 어업인들을 조합원으로 하는 A수산업협동조합(이하 ‘A조합’이라 함)은 조합원들이 포획ㆍ채취한 수산물의 판매를 위탁받아 판매하는 B수산물위탁판매장(이하 ‘B위탁판매장’이라 함)을 운영하여 왔다. 한편, B위탁판매장 운영에 대해서는 관계 법령에 따라 관할 지역에 대한 독점적 지위가 부여되어 있었으며, A조합은 B위탁판매장 판매액 중 일정비율의 수수료를 지급받아 왔다. 그런데, 한국농어촌공사는「공유수면 관리 및 매립에 관한 법률」에 따라 X군 일대에 대한 공유수면매립면허를 받아 공유수면매립사업을 시행하였고, 해당 매립사업의 시행으로 인하여 사업대상지역에서 어업활동을 하던 A조합의 조합원들은 더 이상 조업을 할 수 없게 되었다. A조합은 위 공유수면매립사업지역 밖에서 운영하던 B위탁판매장에서의 위탁판매사업의 대부분을 중단하였고, 결국에는 B위탁판매장을 폐쇄하기에 이르렀다. 이에 따라 A조합은 공유수면매립사업으로 인한 위탁판매수수료 수입의 감소에 따른 영업손실의 보상을 청구하였으나, 한국농어촌공사는 B위탁판매장이 사업시행지 밖에서 운영되던 시설이었고 「공유수면 관리 및 매립에 관한 법률」상 직접적인 보상 규정이 없음을 이유로 보상의 대상이 아니라고 주장한다. 한국농어촌공사의 주장은 타당한가? (20점) \par
\end{longquestion}
